\documentclass[10pt,a4paper]{amsart}
\usepackage[margin=19mm]{geometry}
\usepackage{amsmath,amssymb,mathtools}
\usepackage[scr=rsfs,cal=euler]{mathalfa}
\usepackage{xfrac}
\usepackage[unicode,hidelinks,bookmarks=false]{hyperref}
\newcommand{\ie}{i.e.\ }
\newcommand{\cf}{cf.\ }
\newcommand{\resp}{respectively\ }
\newcommand{\Subsection}{\subsection}
\providecommand{\nobreakdash}{\nobreak\hbox{-}}
\setlength{\emergencystretch}{2em}
\multlinegap0pt
\usepackage{xcolor}
\usepackage{amssymb}
\usepackage{dsfont}
\newtheorem{lemma}{Lemma}
\newtheorem{theorem}{Theorem}
\newcommand{\End}{\op{End}}
\newcommand{\Mm}{\enm{\cal{M}}}
\newcommand{\Oo}{\enm{\cal{O}}}
\newcommand{\PP}{\enm{\mathbb{P}}}
\newcommand{\Uu}{\enm{\cal{U}}}
\newcommand{\cal}[1]{\mathcal{#1}}
\newcommand{\enm}[1]{\ensuremath{#1}}          %
\newcommand{\op}[1]{\operatorname{#1}}
\renewcommand{\to}[1][]{\xrightarrow{\ #1\ }}
\title{On the meagerness of the set of irregular bundles on Hopf surfaces}
\author{Edoardo Ballico, Elizabeth Gasparim}
\date{Source: arXiv:2604.25238v1 (CC BY 4.0)}
\begin{document}
\maketitle

\noindent\emph{Source and licence.} On the meagerness of the set of irregular bundles on Hopf surfaces. Version arXiv:2604.25238v1 (CC BY 4.0), distributed by arXiv under the Creative Commons Attribution 4.0 International licence, http://creativecommons.org/licenses/by/4.0/, as recorded in the arXiv licence metadata for this exact version and shown on the abstract page https://arxiv.org/abs/2604.25238v1. Full licence notice: \texttt{licenses/arxiv-2604.25238-CC-BY-4.0.txt}.\par\smallskip

\noindent\textit{Editorial scope.} Counted unit: the article's theorem i1 (source0.tex lines 361--364). The article's complete original proof of it is delivered with it (source0.tex lines 367--381, one \texttt{proof} environment of 15 lines). No mathematics is written, repaired or re-derived: the statement and the proof are the article's own lines, in the article's own notation, and no retained line is substituted or re-flowed.  Retained uncounted as the source-local context the proof needs: lines 344--346.  Nothing was omitted from any retained span, so the package carries no omission record.  No retained line contains a citation, so the wrapper carries no bibliography and no bibliography entry.  No retained line contains a figure environment, an image-inclusion command or drawing code, so no image is delivered and the package declares no asset dependency.  Editorial formatting only, in addition: an AMS \texttt{amsart} preamble that restates the article's own declarations and macros used by the retained lines, namely the environments \texttt{lemma}, \texttt{theorem} on separate counters, together with the standard packages those lines need.  The retained environments print in this document as Lemma 1, Theorem 1 in order of appearance, whereas the article prints the same items with the numbering of its own full document.  The complete native source of the article as distributed in the arXiv e-print for this exact version is bundled unchanged as \texttt{sources/199287/source0.tex}.
\begin{lemma}\label{i3}
Let $F$ be a rank $r\ge 2$ vector bundle on an elliptic curve $T$ such that $\det(F)\cong \Oo_T$. By Atiyah's classification $F$ is a direct sum of semistable indecomposable ones, $h^0(T,\End(F)) \ge r$ and equality holds if  $F$ is regular. If $F$ is not regular, then $h^0(T,\End(F)) \ge r+2$.
\end{lemma}

\begin{theorem}\label{i1}
Let $\Uu$ be a connected component of $\Mm_{r,c}$. Then the set of all $E\in \Uu$ which are not regular, i.e. either with jumps or irregular, is meager in $\Uu$
(and hence the set of all $E\in \Uu$ which are regular is dense in every open subset of $\Uu$.
\end{theorem}

\begin{proof}
Let $A$ be the set of all $E\in \Uu$ which are not regular, i.e. such that there is $x\in \PP^1$ for which $E_{|F_x}$ is not regular. For each $x\in \PP^1$, let  $A_x$ denote the set of all $E\in \Uu$ such that $E_{|T_x}$ is not regular.
Since $\dim \PP^1 =1$, it is sufficient to prove that each $A_x$ is locally, around each of its points,
 a closed analytic subset of $\Uu$ of codimension at least
$2$. Fix $x\in \PP^1$ and $E\in A_x$. Set $F:= E_{|T_x}$. Consider the exact sequence
\begin{equation}\label{eqd1}
0\to \End(E)(-T_x)\to \End(E)\to \End(F)\to 0 \text{.}
\end{equation}
Duality gives $$h^2(X,\End(E)(-T_x)) =h^0(X,\End(E)(+T_x)\otimes \omega_X) = h^0(X,\End(E)(-T_x)).$$ 
Since $E$ is simple, each element of $H^0(\End(E))$ is induced by the multiplication by a constant. Therefore, $h^0(X,\End(E)(-T_x))=0$.

Thus, the linear map $\rho\colon H^1(X,\End(E))\to H^1(T_x,\End(F))$ is surjective, and an open neighborhood
 of $0\in H^1(X,\End(E))$ is biholomorphic to an open neighborhood of $E$ in $\Mm_{r,c}$. An open neighborhood of $0$ in $H^1(T_x,\End(F))$ is a versal deformation space for $F$. Since $\rho$ is surjective, Lemma \ref{i3} gives that 
$A_x$ has locally at $E$ at least codimension $2$ in $\Uu$.
\end{proof}

\end{document}
